\begin{lemma}
\label{lemma-induction-principle}
Let $S$ be a scheme. Let $X$ be a quasi-compact and quasi-separated
algebraic space over $S$. Let $P$ be a property of the quasi-compact
and quasi-separated objects of $X_{spaces, \etale}$. Assume that
\begin{enumerate}
\item $P$ holds for every affine object of $X_{spaces, \etale}$,
\item for every elementary distinguished square $(U \subset W, f : V \to W)$
such that
\begin{enumerate}
\item $W$ is a quasi-compact and quasi-separated object of
$X_{spaces, \etale}$,
\item $U$ is quasi-compact,
\item $V$ is affine, and
\item $P$ holds for $U$, $V$, and $U \times_W V$,
\end{enumerate}
then $P$ holds for $W$.
\end{enumerate}
Then $P$ holds for every quasi-compact and quasi-separated object
of $X_{spaces, \etale}$ and in particular for $X$.
\end{lemma}

\begin{proof}
We first claim that $P$ holds for every representable
quasi-compact and quasi-separated object of $X_{spaces, \etale}$.
Namely, suppose that $U \to X$ is \'etale and $U$ is a
quasi-compact and quasi-separated scheme. By assumption (1)
property $P$ holds for every affine open of $U$. Moreover, if
$W, V \subset U$ are quasi-compact open with $V$ affine and $P$ holds for 
$W$, $V$, and $W \cap V$, then $P$ holds for $W \cup V$ by (2)
(as the pair $(W \subset W \cup V, V \to W \cup V)$ is an elementary
distinguished square). Thus $P$ holds for $U$ by the induction
principle for schemes, see
Cohomology of Schemes, Lemma \ref{coherent-lemma-induction-principle}.

\medskip\noindent
To finish the proof it suffices to prove $P$ holds for $X$
(because we can simply replace $X$ by any quasi-compact and quasi-separated
object of $X_{spaces, \etale}$ we want to prove the result for).
We will use the filtration
$$
\emptyset = U_{n + 1} \subset
U_n \subset U_{n - 1} \subset \ldots \subset U_1 = X
$$
and the morphisms $f_p : V_p \to U_p$ of
Decent Spaces, Lemma
\ref{decent-spaces-lemma-filter-quasi-compact-quasi-separated}.
We will prove that $P$ holds for $U_p$ by descending induction on $p$.
Note that $P$ holds for $U_{n + 1}$ by (1)
as an empty algebraic space is affine. Assume $P$ holds for $U_{p + 1}$.
Note that $(U_{p + 1} \subset U_p, f_p : V_p \to U_p)$ is an elementary
distinguished square, but (2) may not apply as $V_p$ may not be affine.
However, as $V_p$ is a quasi-compact scheme we may choose a finite affine open
covering $V_p = V_{p, 1} \cup \ldots \cup V_{p, m}$.
Set $W_{p, 0} = U_{p + 1}$ and
$$
W_{p, i} = U_{p + 1} \cup f_p(V_{p, 1} \cup \ldots \cup V_{p, i})
$$
for $i = 1, \ldots, m$. These are quasi-compact open subspaces of $X$.
Then we have
$$
U_{p + 1} = W_{p, 0} \subset
W_{p, 1} \subset \ldots \subset
W_{p, m} = U_p
$$
and the pairs
$$
(W_{p, 0} \subset W_{p, 1}, f_p|_{V_{p, 1}}),
(W_{p, 1} \subset W_{p, 2}, f_p|_{V_{p, 2}}),\ldots,
(W_{p, m - 1} \subset W_{p, m}, f_p|_{V_{p, m}})
$$
are elementary distinguished squares by
Lemma \ref{lemma-make-more-elementary-distinguished-squares}.
Note that $P$ holds for each $V_{p, 1}$ (as affine schemes) and for
$W_{p, i} \times_{W_{p, i + 1}} V_{p, i + 1}$ as this is a quasi-compact
open of $V_{p, i + 1}$ and hence $P$ holds for it by the first paragraph
of this proof. Thus (2) applies to each of these and we inductively
conclude $P$ holds for $W_{p, 1}, \ldots, W_{p, m} = U_p$.
\end{proof}